
\section{De la fuente a la línea del programa}
Cada línea conserva fuente, identificador, revisión, efectividad, acción y regla de intervalo. Su origen determina qué vía de cambio necesita. Una limitación, una AD y una tarea de otro origen no se mezclan en una lista sin procedencia. Un umbral inicial y una repetición pueden tener referencias distintas: alcanzar ciclos totales no equivale a contar desde la última ejecución.

Si se combinan horas y calendario, lo primero, se conservan ambos contadores. La fecha de emisión de un documento no se usa automáticamente como inicio: el disparador procede del requisito. Una revisión puede añadir tareas, cambiar efectividad o conceder crédito condicionado.
\section{Aeronave, motor y pieza tienen historias propias}
El motor puede acumular valores distintos de la aeronave. Una pieza trasladada conserva historial y definición de consumo de vida. Ciclo de vuelo no se transforma automáticamente en ciclo de pieza si sus datos definen otro evento o ponderación. La hoja incluye unidad, definición y procedencia de cada contador.

Límite menos acumulado solo tiene sentido si representan la misma magnitud. Desconocido no equivale a cero. Trasladar una pieza no reinicia total; una revisión general no repone automáticamente toda vida limitada. Las reglas particulares requieren datos y trazabilidad de eventos.
\section{Actualizar sin borrar la justificación}
Se comparan requisitos antiguos y nuevos, unidades afectadas y transición. Después se actualizan programa, planificación y registros con rastro de revisión. Sobrescribir una lista puede borrar por qué una tarea tenía un vencimiento concreto.

La fiabilidad ayuda a evaluar propuestas, pero no elimina obligaciones. Una mejora global no demuestra menor riesgo del modo específico que controla la tarea. Se identifica qué puede cambiarse, con qué evidencia y por qué vía.
\repaso{Quedan 10 h y 12 días para dos límites de la misma tarea. ¿Basta anotar 12 días?}{La utilización puede agotar primero las horas. Se conservan ambos márgenes y previsión de uso, sin modificar el requisito.}
